% en-tete.tex — quelques réglages de mise en page pour les PDF (ajoutés par pandoc dans l'en-tête LaTeX)
\IfFontExistsTF{Charter}{\setmainfont{Charter}}{}
\IfFontExistsTF{Avenir Next}{\setsansfont{Avenir Next}}{}
\usepackage{xcolor}
\usepackage{float}
\floatplacement{figure}{H}                 % chaque figure reste à l'endroit où elle est écrite
\usepackage{ragged2e}
\AtBeginDocument{\RaggedRight}             % texte aligné à gauche (TinyTeX n'a pas la césure du français)
\usepackage{titlesec}
\titleformat*{\section}{\Large\bfseries\sffamily\color[HTML]{1B4F72}}
\titleformat*{\subsection}{\large\bfseries\sffamily\color[HTML]{1B4F72}}
\usepackage{caption}
\captionsetup{font=small,labelfont=bf}
% la table des matières commence sur une nouvelle page, et le texte aussi
\let\anciennetdm\tableofcontents
\renewcommand{\tableofcontents}{\clearpage\anciennetdm\clearpage}
\linespread{1.08}                         % interligne un peu plus aéré
\usepackage{newunicodechar}
\newunicodechar{→}{\ensuremath{\rightarrow}}  % la police Charter n'a pas de flèche
% page de titre : titre plus grand, centré dans le haut de la page
\makeatletter
\renewcommand{\@maketitle}{%
  \null\vspace{2.5cm}%
  \begin{center}%
    {\sffamily\bfseries\huge\color[HTML]{1B4F72}\@title\par}%
    \vspace{2em}%
    {\large\lineskip .6em\begin{tabular}[t]{c}\@author\end{tabular}\par}%
    \vspace{1.5em}%
    {\normalsize\@date\par}%
  \end{center}%
  \vspace{1.5cm}}
\makeatother
